\chapter{Fire Safety in Your Home}

\section*{Definition and Overview}
Home fire safety is the set of practices, equipment, and habits that prevent house fires and protect people when one occurs. Most house fires are preventable, and most fire deaths happen in the home, often at night while people are asleep. In Australia, fire safety starts with working smoke alarms, a practised escape plan, safe cooking and heating habits, and correct use and storage of electrical items, candles, and flammable materials. Fire and Rescue services provide free home safety advice, and every household should know what to do in the first minutes of a fire, because that is when escape is possible. This chapter covers the essentials of preventing fires and surviving them.

\section*{Epidemiology}
Fires in Australian homes cause hundreds of deaths and thousands of injuries each decade. Around half of all house fire deaths occur in homes without a working smoke alarm, and most fatal fires happen at night. Cooking is the leading cause of house fires, followed by electrical faults, heaters, and open flames; smoke alarms, cigarette smoking, and careless heating are common factors in fatal fires. Children, older adults, and people with disability or reduced mobility are at highest risk of death and injury. Working smoke alarms reduce the risk of death in a house fire by around half, making them the single most effective safety measure.

\section*{Aetiology and Pathophysiology}
\begin{itemize}
    \item \textbf{Cooking:} unattended cooking, oil fires, and cooking while affected by alcohol or fatigue are leading causes
    \item \textbf{Electrical faults:} overloaded power boards, damaged cords, faulty appliances, and inappropriate use of power leads
    \item \textbf{Heaters:} portable heaters placed too close to bedding, curtains, or clothing, and heaters left on overnight
    \item \textbf{Smoking:} discarded cigarettes and smoking in bed, especially with alcohol or medicines involved
    \item \textbf{Candles and open flames:} candles left burning unattended, incense, and matches accessible to children
    \item \textbf{Chemical and cooking oil fires:} kitchen oil igniting, and flammable liquids stored near heat
    \item \textbf{How fires kill:} smoke inhalation is the main cause of death; toxic smoke (carbon monoxide and other gases) and heat incapacitate people quickly, often in minutes
    \item \textbf{How fire spreads:} fires double in size rapidly; the escape time in a home fire is often only 2--3 minutes
    \item \textbf{The role of smoke alarms:} they provide the early warning that makes escape possible
\end{itemize}

\section*{Clinical Features}
Fire injuries and the situations that cause them have recognisable features.
\begin{itemize}
    \item Smoke alarms: photoelectric alarms are required in all Australian homes; they should be installed on each level, in hallways, and near bedrooms
    \item Burns and smoke inhalation are the classic injuries, affecting the skin and airways
    \item Common risk situations: cooking left unattended, overloaded power boards, and heaters near bedding
    \item The early warning signs of fire risk in a home: no working smoke alarms, no escape plan, and fire hazards in daily use
    \item Carbon monoxide poisoning, which is odourless and can occur with poorly maintained heaters
    \item The first signs of a fire: smoke, smell of burning, and the sound of the alarm
\end{itemize}

\section*{Differential Diagnosis}
Assessing fire risk involves checking common hazards.
\begin{itemize}
    \item Working smoke alarms in all required locations, with batteries checked regularly
    \item Cooking practices: never leaving cooking unattended, and keeping combustibles away from the stove
    \item Electrical safety: no overloaded power boards, damaged cords, or faulty appliances
    \item Heating safety: heaters at least one metre from combustibles and never left on while asleep
    \item Smoking safety: smoking outside and safe disposal of butts
    \item Candles and matches stored safely, especially where children are present
    \item The presence of fire extinguishers and fire blankets, and knowledge of how to use them
\end{itemize}

\section*{When to Seek Medical Care}
Anyone with burns, smoke inhalation, or any injury from a fire needs medical assessment. Burns that are larger than a small area, on the face, hands, or joints, or that blister extensively should be treated as emergencies.

\textbf{Call 000 (or local emergency services) for:}
\begin{itemize}
    \item Any fire, even a small one; call 000 and get everyone out first
    \item Burns with blistering, burns to the face, hands, or groin, or burns larger than the palm of the hand
    \item Difficulty breathing, coughing, or a hoarse voice after a fire or smoke exposure
    \item Signs of smoke inhalation: confusion, headache, or loss of consciousness in a fire setting
    \item Anyone trapped or unaccounted for
\end{itemize}

\section*{Diagnosis and Investigations}
\begin{itemize}
    \item \textbf{Home safety audit:} check smoke alarms, escape routes, and fire hazards room by room; Fire and Rescue NSW and other services offer free home safety visits
    \item \textbf{Smoke alarm testing:} alarms should be tested monthly and batteries replaced at least yearly, or the sealed unit replaced every 10 years
    \item \textbf{Escape plan practice:} every household should have two ways out of every room and a meeting point outside
    \item \textbf{Assessment of specific hazards:} electrical, heating, cooking, and smoking practices reviewed and corrected
    \item \textbf{Medical assessment:} for burns and smoke inhalation, including airway and lung examination
    \item \textbf{Carbon monoxide checks:} for households using gas, wood, or solid fuel heaters
\end{itemize}

\section*{Management}
\begin{itemize}
    \item \textbf{Install and maintain smoke alarms:} photoelectric alarms on every level, tested monthly, with batteries replaced yearly and units replaced at 10 years
    \item \textbf{Practise an escape plan:} identify two exits from every room, keep exits clear, and agree on a meeting point; practise the plan with the family, including night-time practice
    \item \textbf{Cook safely:} never leave cooking unattended, keep pot handles turned in, and keep a fire blanket near the kitchen
    \item \textbf{Heat safely:} keep heaters a metre from combustibles, never dry clothes on heaters, and service heaters annually
    \item \textbf{Electrical safety:} use licensed electricians, avoid overloaded power boards, and replace damaged cords
    \item \textbf{Smoke safely:} smoke outdoors, use deep ashtrays, and never smoke in bed
    \item \textbf{Candle safety:} never leave candles burning unattended and keep them away from curtains
    \item \textbf{Know the drill:} if a fire starts, get everyone out, close doors behind you, and call 000 from outside; never go back inside
\end{itemize}

\section*{Complications}
\begin{itemize}
    \item Burns, which can be disfiguring and disabling, and require specialist care
    \item Smoke inhalation and carbon monoxide poisoning, which can be fatal or cause brain injury
    \item Death, most commonly at night and in homes without working smoke alarms
    \item Loss of home, possessions, and pets
    \item Post-traumatic stress and grief after a fire
    \item High financial costs of fire damage and insurance shortfalls
\end{itemize}

\section*{Prognosis and Outcomes}
Most home fires are preventable, and the outcomes of fires are dramatically better when smoke alarms work and escape plans are practised. The majority of people who die in house fires are in homes without a working alarm, so the single most effective action is installing and maintaining alarms. Burns and smoke inhalation are survivable with prompt emergency care, but the best outcome is avoiding the fire altogether. Households that treat fire safety as a routine habit, like locking the door, virtually eliminate the risk of fire death.

\section*{Prevention and Self-Care}
\begin{itemize}
    \item Install photoelectric smoke alarms on every level and near bedrooms, and test them monthly
    \item Replace alarm batteries yearly and units every 10 years
    \item Draw and practise a home escape plan with two ways out of every room
    \item Never leave cooking, candles, or heaters unattended
    \item Keep heaters, lamps, and appliances away from curtains and bedding
    \item Avoid overloaded power boards and damaged electrical cords
    \item Keep matches and lighters out of children's reach
    \item Store flammables outside or safely away from heat sources
    \item Ask your state fire service for a free home safety visit
\end{itemize}

\section*{Sources and Further Reading}
The following reputable sources provide further information for healthcare students and patients seeking reliable, current advice about home fire safety.
\begin{itemize}
    \item Fire and Rescue NSW. \textit{Home fire safety.} https://www.fire.nsw.gov.au/
    \item Country Fire Authority (Victoria). \textit{Home fire safety.} https://www.cfa.vic.gov.au/
    \item Queensland Fire and Emergency Services. \textit{Home safety.} https://www.qfes.qld.gov.au/
    \item Fire and Emergency Services of WA. \textit{Smoke alarms and escape plans.} https://www.dfes.wa.gov.au/
    \item US National Fire Protection Association. \textit{Home fire safety.} https://www.nfpa.org/
    \item World Health Organization. \textit{Burn prevention.} https://www.who.int/
\end{itemize}
